\begin{textsample}{POS Dim 3 – to – Score 22.00 – to/to\_inf\_20.txt}  \label{ex:f3_pos_109}
\textbf{Criança} em foco : planejamento na educação \textbf{infantil} Planejar na Educação \textbf{Infantil} é voltar -se para o ato da reflexão contínua sobre a prática educativa, parte da compreensão de quem são as \textbf{crianças}, grupos \textbf{etários} atendidos, gostos e preferências, modo de vida e de onde vieram, bagagens culturais, seus interesses e de sua comunidade.
Nesse sentido, as DCNEI ( 2009 ) defendem que o planejamento norteia o currículo como um conjunto de práticas que buscam articular as experiências e os saberes das \textbf{crianças} com os conhecimentos que fazem parte do patrimônio cultural.
Para tanto, é importante que o planejamento seja visto como uma oportunidade de autoria criativa do trabalho pedagógico e cabe ao professor considerar as experiências e os conhecimentos de mundo das \textbf{crianças} e, a partir das referências, compreender e garantir práticas contextualizadas e narrativas permeadas pelas interações e \textbf{brincadeiras}.
Com ações organizadas pelo professor e auto-organizadas pelas \textbf{crianças}, atribuem -se diferentes sentidos para a intencionalidade pedagógica, pois são elas, as \textbf{crianças}, que constroem suas próprias aprendizagens.
O professor oferece contextos com intencionalidades e provocações que permitem à \textbf{criança} conviver, \textbf{brincar}, participar, explorar, experimentar e conhecer.
Planejar situações que provocam a reflexão e ação da \textbf{criança}, envolvendo o tempo, os espaços da \textbf{instituição}, as diferentes linguagens, os espaços \textbf{lúdicos} ( sessões, cantos de interesses, ateliês, territórios, oficinas, entre outros ), garantindo os direitos de aprendizagem articulados aos campos de experiências.
Portanto, o professor deve desenvolver um planejamento flexível, que possa ser direcionado conforme o interesse e expectativas das \textbf{crianças}, segundo preconizam as DCNEI ( 2009 ) : “ em relação a qualquer experiência de aprendizagem que seja trabalhada pelas \textbf{crianças}, devem ser abolidos os procedimentos que não reconhecem a atividade criadora e o protagonismo da \textbf{criança} \textbf{pequena}, que promovam atividades mecânicas e não significativas para as \textbf{crianças} ”.
Arte de Eloah Holanda Nascimento, 05 anos O desafio é promover propostas significativas para as \textbf{crianças} que irão ressignificar essas experiências, promovendo situações para além do \textbf{cuidar} e educar.
Os professores organizam ações básicas para o exercício da profissão docente : a observação, o \textbf{registro}, a problematização e a verificação da proposta, ou seja, a documentação pedagógica.
Tais ações, quando incorporadas como atividade docente, constituem -se em preciosos instrumentos que potencializam o trabalho contínuo de planejamento e avaliação.
Essas práticas fazem com que o planejamento seja sinônimo de uma atividade sempre inovadora e diferente a cada ano, de acordo com as diferentes turmas de \textbf{crianças}.
Para tanto, segundo Barbosa ( 2006 ), o planejamento das propostas feitas às \textbf{crianças} precisa atender a alguns critérios básicos : Ao planejar, o professor analisa o perfil da turma, a proposta ( projetos ) desejada pelas \textbf{crianças} ( interesses ), busca contemplar as experiências no contexto \textbf{diário}, semanal, mensal, bimestral.
Situa -se no ambiente escolar, ou seja, nos espaços a serem apropriados pelas \textbf{crianças}, áreas internas e externas, cantos de interesses, entre outros.
Finalmente, ele toma decisões com foco nas linguagens, espaços, tempos, materiais, nos campos de experiências, nos direitos de aprendizagem e na busca por qualidade das interações.
A seguir, apresentam -se \textbf{sugestões} de uso de diferentes instrumentos do planejamento do trabalho pedagógico, que podem ajudar a organizar as propostas no tempo, partindo de uma unidade maior, que é o plano anual, até chegar à unidade \textbf{menor}, o planejamento \textbf{diário}.
Com base na proposta pedagógica da \textbf{instituição} e considerando as aprendizagens selecionadas para o grupo de \textbf{crianças}, os professores precisam organizar \textbf{rotinas} \textbf{diárias}, considerar diferentes modalidades organizativas que levem em conta sua frequência e encadeamento.
Nesse sentido, uma \textbf{sugestão} de organização do tempo didático das \textbf{instituições} de Educação \textbf{Infantil} tem sido : a ) Níveis do planejamento Plano anual O plano anual é uma ideia geral e é organizado globalmente, considerando a experiência anterior da \textbf{criança} e o que ela ainda precisa aprender.
Nesse plano estão previstas as experiências organizadas em torno das diferentes linguagens na construção dos objetos de conhecimento.
Mas para distribuir tais experiências ao longo do ano é necessário que o professor reflita sobre como vai propor : em forma de projetos, atividades permanentes, atividades ocasionais, entre outras.
Planejamento \textbf{diário} É um documento específico que orienta o trabalho com as \textbf{crianças} em um determinado \textbf{dia}.
Esse instrumento é útil ao professor na organização de suas ideias e hipótese de trabalho em relação ao que ele espera garantir ao grupo de \textbf{crianças}.
Ao pensar antecipadamente o que vai propor às \textbf{crianças}, o professor pode acionar o que sabe sobre seu grupo, explicitar com mais clareza seus objetivos de trabalho e refletir sobre as tantas variáveis que podem interferir nas experiências das \textbf{crianças}.
Para tanto, precisa estar \textbf{apoiado} na jornada do grupo, articulando o fluxo de \textbf{crianças} de outros grupos nos mesmos espaços e nas \textbf{rotinas} de serviços de apoio.
Têm objetivos claros com intencionalidade pedagógica, divisão de etapas que podem se relacionar com temáticas divididas em grupos de situações de aprendizagens permeadas pelas interações e \textbf{brincadeiras}, garantindo os direitos de aprendizagem articulados com campos de experiências. b ) Projetos Os projetos da escola necessitam dialogar com as Diretrizes Curriculares Nacionais, Base Nacional Comum Curricular, com o Currículo do Território e a Proposta Pedagógica da Unidade.
A unidade pode desenvolver o seu Projeto institucional conforme as características culturais da comunidade, pontos de interesses e necessidades educativas.
É importante que as \textbf{crianças} e professores tenham espaço para dialogar com a proposta, que necessita ser flexível à mudança. c ) O Projeto da turma : construção ou investigação Ao trabalhar com projeto, o professor da Educação \textbf{Infantil} necessita assumir a postura de observador das relações entre as \textbf{crianças} e entre elas e os \textbf{adultos}, conhecer seus pontos de interesses, seus questionamentos sobre o mundo, assumir a postura de uma \textbf{escuta} sensível.
Nesse sentido, Brasil ( 2016 ) afirma que essa é uma postura fundamental para dar visibilidade às \textbf{crianças} e às suas manifestações.
O professor, ao se colocar numa posição de \textbf{escuta}, oferece às \textbf{crianças} a possibilidade de se colocarem de modo criativo na realidade.
É o reconhecimento da \textbf{criança} como alguém que intervém e transforma a realidade à medida que é por ela transformado.
Para elaborar o projeto, o professor necessita refletir : - O tema do projeto vem ao encontro dos interesses e \textbf{desejos} das \textbf{crianças} da turma, ou é o professor que deseja trabalhar o tema com foco na necessidade das \textbf{crianças}? - Quais conhecimentos as \textbf{crianças} têm sobre o tema/proposta? - Quais as possibilidades das \textbf{crianças} terem contato com a proposta? - De que forma a família pode ser estimulada a participar e se envolver com a proposta? - Quais questionamentos devem ser feitos para as \textbf{crianças} e suas famílias? - Como os campos de experiências e os direitos podem se comunicar dentro da proposta? - Quais são os \textbf{desejos} e anseios das \textbf{crianças} voltados para a proposta? - Como os espaços, dentro e fora da unidade, podem ser utilizados? - O projeto terá o envolvimento de artistas, artesãos, comunidade, grupos intergeracionais, étnicos, pais, profissionais, entre outros? - Quais recursos e materiais serão utilizados na proposta? - Qual o tempo destinado para a execução da proposta?
Deste modo, o planejamento dialoga com o pensamento de Kramer : “ ver e \textbf{ouvir} são cruciais para que se possam compreender \textbf{gestos}, discursos e ações, para descobrirmos o que as \textbf{crianças} já sabem e como constroem significados para o mundo ” ( KRAMER, 2013, p. 48 ).
O que já sabemos Preparação Levantamento de questões e saberes das \textbf{crianças}.
O que queremos saber?
Problematização Definir questões que suscitaram a \textbf{curiosidade} das \textbf{crianças}.
Como vamos saber?
Pesquisa e experimentação Definir caminhos para investigação.
O que aprendemos?
Resultados Das experiências e descobertas.
Confrontar com as respostas iniciais.
Como o próprio nome define, são situações que acontecem todos os \textbf{dias}, estão intrinsecamente ligadas à \textbf{rotina}, leitura ou contação de histórias, \textbf{brincadeiras} livres e dirigidas, desenho, interações com a água, cantos de interesses ; tais situações constituem oportunidades para o desenvolvimento \textbf{infantil}.
Permitem trabalhar com as \textbf{crianças} um contexto/tema que se considera valioso, mesmo não tendo correspondência com o que está planejado para o momento.
A verificação da proposta é o ato de avaliar as experiências oferecidas, se elas são permeadas pelos campos de experiências e os direitos de aprendizagem.
Uma proposta completa garante experiências significativas para as \textbf{crianças}, enriquece repertório, promove a autonomia, favorece o imaginário, as interações e as \textbf{brincadeiras}.
As modalidades organizativas não devem ser trabalhadas isoladamente, mas em interação, favorecendo

% matched lemmas: adulto, apoiar, brincadeira, brincar, criança, cuidar, curiosidade, desejo, dia, diário, escuta, etário, gesto, infantil, instituição, lúdico, ouvir, pequeno, registro, rotina, sugestão
\end{textsample}
